\documentclass[border=4pt]{standalone}
\usepackage[siunitx]{circuitikz}

\begin{document}
\begin{circuitikz}[european resistors]
  % One source, one load. Fix the reference direction of I (clockwise) and
  % the polarity of V_R (+ where I enters R) BEFORE reading any signed result.
  \draw (0,0) node[ground] {}
    to[V, l_=$V_{S}$] (0,3)            % source; + at the top rail
    to[short, i>=$I$] (2,3)            % current reference arrow, clockwise
    to[R, l=$R$] (2,0)                 % load
    -- (0,0);                          % return rail to the reference node

  % Explicit connection dots where the test-point branches make 3-conductor
  % junctions at the load terminals.
  \node[circ] at (2,3) {};
  \node[circ] at (2,0) {};

  % Source polarity marked explicitly at both terminals.
  \node[font=\small] at (-0.34,2.55) {$+$};
  \node[font=\small] at (-0.34,0.45) {$-$};

  % Voltage across R marked by hand: + where I enters, - where it leaves,
  % with V_R in clear space between the source and R.
  \node[font=\small] at (1.6,2.55) {$+$};
  \node[font=\small] at (1.6,0.45) {$-$};
  \node[font=\small] at (1.32,1.5) {$V_{R}$};

  % Test points brought out to clear space; labels sit off the wires.
  \draw (2,3) -- (3.5,3)
    node[ocirc, label={[font=\small]right:{TP\_A}}] {};
  \draw (2,0) -- (3.5,0)
    node[ocirc, label={[font=\small]right:{TP\_B ($0$~V)}}] {};
\end{circuitikz}
\end{document}
